\par\Needspace{6\baselineskip}\medskip
\begin{center}\hypertarget{native-222-article160}{}\pdfbookmark[2]{160.}{article-160}160.\end{center}
In expositione methodi, cui disquisitiones praecedentes dicatae fuerunt,
in quosdam casus speciales incidimus, ubi applicationem non patitur, saltem
non in forma ea, in qua a nobis exhibita est. Hunc defectum locum habere
vidimus:

\SourceLeaf{223}\noindent\emph{primo}, quoties aliquis trium locorum geocentricorum vel cum loco respondente
heliocentrico terrae, vel cum puncto opposito coincidit (casus posterior
manifesto tunc tantum occurrere potest, ubi corpus coeleste inter Solem
et terram transiit),

\noindent\emph{secundo}, quoties locus geocentricus primus corporis coelestis cum tertio
coincidit,

\noindent\emph{tertio}, quoties omnes tres loci geocentrici una cum loco heliocentrico terrae
secundo in eodem circulo maximo siti sunt.

In casu primo situs alicuius circulorum maximorum $AB$, $A'B'$, $A''B''$
indeterminatus manebit, in secundo atque tertio situs puncti $B^*$. In hisce
itaque casibus methodi supra expositae, per quas, si quantitates $P$, $Q$ tamquam
cognitae spectantur, e locis geocentricis heliocentricos determinare docuimus,
vim suam perdunt: attamen discrimen essentiale hic notandum est, scilicet in
casu primo hic defectus soli methodo attribuendus erit, in casu secundo et
tertio autem ipsius problematis naturae; in casu primo itaque ista determinatio
utique effici poterit, si modo methodus apte varietur, in secundo et
tertio autem absolute impossibilis erit, locique heliocentrici indeterminati
manebunt. Haud pigebit, hasce relationes paucis evolvere: omnia vero, quae ad
hoc argumentum pertinent, exhaurire eo minus e re esset, quod in omnibus
his casibus specialibus orbitae determinatio exacta impossibilis est, ubi a
levissimis observationum erroribus enormiter afficeretur. Idem defectus etiamnum
valebit, quoties observationes haud quidem exacte, attamen proxime ad
aliquem horum casuum referuntur: quamobrem in eligendis observationibus
huc respiciendum, probeque cavendum est, ne adhibeatur ullus locus, ubi
corpus coeleste simul in viciniis nodi atque oppositionis vel coniunctionis
versatur, neque observationes tales, ubi corpus coeleste in ultima ad eundem
locum geocentricum proxime rediit, quem in prima occupaverat, neque demum
tales, ubi circulus maximus a loco heliocentrico terrae medio ad locum
geocentricum medium corporis coelestis ductus angulum acutissimum cum
directione motus geocentrici format, atque locum primum et tertium quasi
stringit.
